\section{Results}
\subsection{Volatile memristive device}

The volatile synapse used in this work is a resistive switching C\,/\,HfO$\mathrm{_2}$\,/\,Ag memristive device (see fabrication details in \textit{Materials and Methods}). After the device fabrication, the device is in its high-resistive state and can be switched to the \gls{lrs} by applying a quasi-static voltage sweep between 0\,V and 1.5\,V and back. Contrary to many filamentary devices, the proposed device does not require an electroforming operation (Fig.~\ref{fig:Supfig_01}a), thus simplifying the electrical operation within the neuromorphic circuit. The maximum current flowing through the device (i.e., current compliance, I$\mathrm{_{CC}}$) is limited by applying a voltage to the gate of an NMOS transistor whose drain is connected to the bottom electrode of the memristive device (Fig.~\ref{fig:fig2}A, inset). The switching to the \gls{lrs} occurs because, as the voltage exceeds a given \textit{threshold voltage} V$\mathrm{_{T}}$, Ag ions migrate across the oxide layer and form the \gls{cf}, thus bringing the device in a \gls{lrs} (\textit{set operation}, Fig.~\ref{fig:fig2}A) \cite{Wang2016, Li2020, wang2019TED-II}. However, the filament self-sustains only in presence of a voltage higher than a critical voltage referred to as \textit{hold voltage} V$\mathrm{_{H}}$. Below this value, the rediffusion of the Ag atoms in the dielectric layer causes the self-disruption of the filament and, as a consequence, the device transition to a \gls{hrs} (\textit{spontaneous recovery}, Fig.~\ref{fig:fig2}A). These operations are fairly reproducible (Fig.~\ref{fig:Supfig_01}b). The two main features that enable the use of the proposed memristive device as a volatile synapse in working memory tasks are tunable volatility in biologically relevant time-scales and controllable switching probability.

The retention time is shown in Fig.~\ref{fig:fig2}B and the median retention time increases exponentially with the current compliance, as illustrated in Fig.~\ref{fig:fig2}C. The device can be tuned to be in \gls{lrs} for a time ranging from ms to seconds depending on I$\mathrm{_{CC}}$, which is the time-scale of interest for our application. The intrinsic stochasticity of the filament formation at a microscopic level originates high variability in the distributions of the retention times \cite{Chekol2021}.

The controllability of the switching probability (P$\mathrm{_{ON}}$) of the device can instead be achieved by either changing the width or the amplitude of the voltage pulse applied to the device, as shown in Fig.~\ref{fig:fig2}D.  
It is worth noticing that the device operates with voltages $<$3\,V, which are compatible with standard \gls{cmos} technology and therefore suitable for integration with neuromorphic circuits. Fig.~\ref{fig:fig2}E shows the stochastic behavior of the device. Due to its stochastic nature, the threshold voltage of the device is slightly variable, which results in an electrically tunable switching probability: The higher the pulse voltage, the higher the switching probability \cite{covi2021TED-I}. Therefore, the desired switching probability of the device can be tuned by setting a suitable voltage amplitude. Moreover, when stimulated by a burst of identical pulses (Fig.~\ref{fig:Supfig_02}), the switching probability of the device increases with the number of pulses, as demonstrated in Fig.~\ref{fig:fig2}F and \ref{fig:Supfig_03bis}. %This behavior is dependent on the stimulation frequency: If the burst occurs at low frequencies (e.g., $<$10\,Hz), there is no evident dependence of the switching probability on the number of pulses \erika{I remember a discussion about it, but we need to retrieve the data and add them to the SI}.
This effect can therefore be exploited to accelerate the storing or training phase in a network.

\begin{figure}[ht]
    \centering
    \includegraphics[width=0.95\textwidth]{Figures/Fig2.pdf}
    \caption{Ag-based volatile memristive device characterization. (a) Sketch of the one-transistor\,/\,one-resistor (1T1R) RRAM device together with its working principle. The memristive device (1R) is based on a W\,/\,C\,/\,10\,nm HfO$\mathrm{_2}$\,/\,Ag stack. The RRAM shows a volatile behavior, i.e., a set operation together with a spontaneous switch off. (b) Time characterization of the retention of the filament. After a 5\,V amplitude triangular pulse to switch  the cell on with an I$\mathrm{_{CC}}$\,=\,20\,\textmu A, a constant reading voltage of -150\,mV is applied to monitor the retention. (c) Retention time distributions at different compliance currents. The median value of the retention time increases with the compliance current. Inset: applied programming pulse. (d) Switching probability of the device for a single pulse as a function of the amplitude and the pulse duration of the programming pulse. Shorter pulses required higher voltage amplitudes to switch ON. (e) Probability of switching the RRAM to the \gls{lrs} depending on the number of pulses and their amplitude (1\,ms pulses). The circles are the experimental data while the solid lines is the fitting. (f) Effect of the number of pulses and the voltage amplitude (1.6\,V) on the switching of the device. The switching of the device is stochastic. Considering a group (burst) of pulses, the probability that the device is in the \gls{lrs} inside the group increases.}
    \label{fig:fig2}
\end{figure}

\subsection{Working memory}
\subsubsection{Store and recall of features}

The volatile behavior of the memristive device is used in a small-scale experiment of \gls{wm} to store and recall features, as depicted in Fig.~\ref{fig:fig3}A. In our example, our network consists of 5 volatile synapses all connected to the same neuron (schematic in Fig.~\ref{fig:fig3}B). The current flowing to the neuron is the sum of the currents flowing through the stimulated devices and it depends on the resistive state of each device. The neuron is trained to recognize a color within a color stream. Each color is encoded by stimulating a different combination of three devices as in Fig.~\ref{fig:fig3}C. As a result, for each color only three devices out of five are switched on. At first, the target color is stored in the network by repeatedly stimulating the network with the combination of pulses that encodes the desired color. As a consequence, the corresponding stimulated devices will be switched on and set to their \gls{lrs}. Afterwards the network is fed with a stream of random stimuli, as shown in Fig.~\ref{fig:fig3}D (see also Fig.~\ref{fig:Supfig_03}). When the stored item is presented to the network, all the stimulated devices are in \gls{lrs} and therefore the current received by the neuron overcomes a threshold, thus triggering the firing of the neuron. The current threshold is set according to the current distributions shown in Fig.~\ref{fig:Supfig_04}. After about 1\,s without stimulation, all the devices switch off and the pattern is forgotten, as shown in Fig.~\ref{fig:Supfig_03}. The correlation plot in Fig.~\ref{fig:fig3}E shows the difference between the expected and the measured current when a pattern is presented. The main contribution to the current is given by the synapses that are stimulated and in \gls{lrs}, therefore three different current levels can be expected. It is shown that we can reach more than 90\% accuracy, defined as the percentage of correct classifications (stored / non-stored pattern) during a test sequence. %\erika{@Saverio, can you confirm it, please?} %To deeply understand the behavior of the system, different parameters have been tested, in terms of spike rate and switching probability. The spike rate (10, 20 and 50\,Hz) is mainly related to the retention time, since if the rate is too slow, the retention times of the devices make them coming from the LRS to the HRS and so to forget the stored information. The switching probability (PON = 5\%, 10\%, 20\%) instead impacts on the robustness of the system to keep the information without changing it. The higher the switching probability and the faster the system become corrupted and the stored item changes. By properly tuning these two parameters is it possible to maximize the efficiency, the robustness and the retention of the system.

\begin{figure}[ht!]
    \centering
    \includegraphics[width=0.95\textwidth]{Figures/Fig3.pdf}
    \caption{\Gls{wm} store and recall experiment. (a) High-level sketch of the working memory. (b) Schematic of the \gls{wm} implementation: 5 volatile 1T1R memristive devices are arranged in parallel configuration. The gate is chosen to set I$\mathrm{_{CC}}$\,=\,17\,\textmu A, that corresponds to a retention time of 28\,ms. (c) Color - pattern encoding. (d) Store and recall experiment. During the store phase, a single pattern is fed to the network. Top colored plot: input stimuli. For ease of visualization, each pattern is colored as the color it encodes. Black dots in the bottom part of the upper plot indicate the stored pattern. Bottom plot: measure current fed to the post-neuron. The current threshold for recognition is indicated as a dashed black horizontal line. The traces are cropped on the x-axes to better highlight the salient events. (e) Correlation plot between the expected and measured currents based on the difference between the presented and the stored pattern. Results obtained from 10 different store and recall experiments with P$\mathrm{_{ON}}$\,=\,5\% and stimulation frequency f$\mathrm{_{stim}}$\,=\,50\,Hz. (f) Accuracy of the system in distinguishing the stored pattern under different stimulation and switching conditions. (g) Average current error, defined as the difference between the measured current and the expected current, during 100 patterns applied for the different conditions. %The P$\mathrm{_{ON}}$\,=\,5\% with f$\mathrm{_{stim}}$\,=\,50\,Hz resulted to be the best operative condition, with a small loss of current only at the end of the experiments.
    }
    \label{fig:fig3}
\end{figure}

The retention time and the switching probability of each synaptic device are electrically tunable through the gate voltage of the transistor (which sets the current compliance) and the voltage amplitude of the pulse applied to the \gls{te} of the device, respectively. These features, together with the stimulation frequency of the network, determine the accuracy of the network in recognizing the color. As shown in Fig.~\ref{fig:fig3}F, for low P$\mathrm{_{ON}}$ (P$\mathrm{_{ON}}$\,=\,5\%), an increase of the spike rate leads to an improvement of the accuracy because the time between two pulses is shorter than the retention time of the device. Instead, higher values of P$\mathrm{_{ON}}$ lead the network to experience a decrease in its accuracy because the probability that a non-active device is switched on is higher. Finally, the intrinsic volatility of the network results in a progressive forgetting of the stored element. Fig.~\ref{fig:fig3}G gives an estimation on how long the network remembers by showing the average mismatch current, defined as the difference between the expected current and the measured current. The two main errors that can occur over time, i.e., the network either fails to recognize the correct color or recognizes the wrong one, depend on the combination of stimulation frequency and P$\mathrm{_{ON}}$. Indeed, for high P$\mathrm{_{ON}}$ and stimulation frequency, devices that were supposed to be in their \gls{hrs} turn to their \gls{lrs}, thus increasing the measured current. The opposite combination, i.e., low P$\mathrm{_{ON}}$ and stimulation frequency, results in the switching off of devices previously in \gls{lrs}, thus decreasing the measured current. The results suggest that indeed a careful selection of the P$\mathrm{_{ON}}$ based on the planned stimulation frequency is advisable to fine tune the \gls{wm} accuracy for a given task, as it will be discussed further in Section~\ref{sec:disc}.
